By Vieta, $\alpha+\beta=2$ and $\alpha\beta=-5$. The denominator product $(\alpha+2)(\beta+2)=3$ is nonzero.

(i) Put $u=\dfrac{\alpha}{\alpha+2}$, $v=\dfrac{\beta}{\beta+2}$. Then $u+v=-2$, $uv=-\dfrac{5}{3}$. If $s_j=u^j+v^j$, then
$s_j=-2s_{j-1}+\dfrac{5}{3}s_{j-2}$, with $s_0=2,s_1=-2$. Successive values are
\[
s_2=\dfrac{22}{3},\quad s_3=-18,\quad s_4=\dfrac{434}{9},\quad
\boxed{s_5=-\dfrac{1138}{9}}.
\]
(ii) Polynomial division gives
\[
\dfrac{x^5}{x+2}=x^4-2x^3+4x^2-8x+16-\dfrac{32}{x+2}.
\]
The original power sums satisfy $p_j=2p_{j-1}+5p_{j-2}$, so $p_1=2,p_2=14,p_3=38,p_4=146$. Also $\dfrac{1}{\alpha+2}+\dfrac{1}{\beta+2}=2$. Therefore the answer is
$146-2(38)+4(14)-8(2)+32-64=\boxed{78}$.
